\begin{tablaObservaciones}{Informe 2}
\grupoObservaciones{Presentación y control transversal}
\observacion[sub=E]{GEN-01}{GEN-01. Las tablas de respuesta ocupan una sección que altera el índice obligatorio de los subdocumentos. (rev., pp. 3)}{Centralizar las respuestas en el anexo del Art. 46 e incorporarlo al índice general del Informe 3. Retirar las secciones de respuesta antepuestas a los capítulos. *(Implementado en D1: verificado y confirmado que las aperturas de S1, S2 y S9 en Informe-3/d1/ carecen de tablas de Art. 46 antepuestas, centralizando el control en este anexo y en el registro inter-áreas responsables).*}{Índice general del Informe 3 y aperturas de S1, S2 y S9 (Informe-3/d1/). Estado: Incorporada en D1}
\observacion[sub=E]{GEN-02}{GEN-02. Los folios reinician por archivo. La revisión exige numeración correlativa del conjunto para la oferta final. (rev., pp. 1–3, 14, 47)}{Asignar el folio global al ensamblar la oferta y conciliar el inicio de cada documento con la hoja resumen de su sobre. Verificar media firma, foliación y firmas de apertura según Art. 40.}{Hojas resumen y pies de página de la oferta final. Estado: Sin cierre acreditado}
\observacion[sub=E]{GEN-03}{GEN-03. S1, S2 y T-6 reiteran cláusulas defensivas. S1 mantiene prosa institucional de folleto. (rev., pp. 3–5, 7–8, 47)}{Se depuró la presentación de capacidades, compromisos y supuestos. S1 retiró misión, visión y valores. El T-6 concentra el cotejo de cartas e instrumentos de recepción en la referencia al Sobre N.º 1. *(Complementado en D1: T-6 f. 6 enlaza expresamente a la subsección de absorción de escala de S1 §1.4).*}{S1 §1.1.1, f. 8, y cuerpo. S2 §§2.1–2.5, desde f. 8. T-6, f. 6 (Informe-3/d1/). Estado: Incorporada}
\observacion[sub=E]{GEN-04}{GEN-04. Las declaraciones de S3, S4 y S6 tienen seis columnas y la revisión las observa por el límite general del cuerpo. (rev., pp. 11, 14, 23, 47)}{Conservar los seis campos expresamente exigidos por el Comunicado 10 §7.2 y distinguir la declaración obligatoria de las tablas analíticas del cuerpo. Aplicar el mismo criterio a las tres declaraciones y conciliarlo con la contraparte si se exige una presentación distinta.}{Declaraciones de uso de IA de S3, S4 y S6. Estado: Sin cierre acreditado}
\observacion[sub=E]{GEN-05}{GEN-05. S6, S8 y S9 declaran uso alto de IA con revisión concentrada. Se solicita especial atención a la simulación de riesgos y al diseño de pruebas. (rev., pp. 23, 31–32)}{Verificar fuentes y decisiones de S6, reproducir el modelo de riesgos a partir del T-15 y verificar los casos de prueba contra sus oráculos. Consignar el uso real y las verificaciones humanas en las declaraciones y el A-6. *(Implementado en D1: emitido el registro formal de auditoría humana Informe-3/d1/REGISTRO\_REVISION\_HUMANA\_GEN05.md documentando la verificación y validación caso a caso de los 120 ensayos de T-17, 13 protocolos de hitos y 29 criterios por Carlos Abarza; actualizadas las declaraciones de IA en S1, S2 y S9).*}{Declaraciones de S1, S2, S9 y REGISTRO\_REVISION\_HUMANA\_GEN05.md (Informe-3/d1/). Pendiente en S6 y S8 por D3/D4. Estado: Incorporada en D1 (Parcial global)}
\observacion[sub=E]{GEN-06}{GEN-06. El T-11 carece de declaración propia y los T-14, T-15 y T-18 remiten a la del capítulo. (rev., pp. 17–18, 27, 47)}{Individualizar la declaración por formulario según la observación, conservar las filas asociadas en el subdocumento y consolidar el registro con A-6. En S7 verificar también la mención expresa del T-15.}{T-11, T-14, T-15, T-18, declaraciones de S4 y S7 y A-6. Estado: Sin cierre acreditado}
\grupoObservaciones{Empresa y comprensión de la necesidad}
\observacion[sub=E]{EMP-01}{EMP-01. La experiencia de mayor escala declarada cubre 120 tractocamiones frente a 374 camiones y 96.000 viajes anuales del caso, sin explicar cómo se compensa la brecha. (rev., pp. 4–5, 39)}{Mantener la escala acreditada de cada referencia. Justificar la capacidad para Curimón con dotación, refuerzos verificables, turnos, red regional y productividad conciliada con el plan, sin atribuir equivalencias entre sensores y camiones. *(Implementado en D1: agregada en S1 §1.4 la subsección «Capacidad operativa y absorción de escala para el contrato» con Tabla 1-capacidad-frentes, justificando absorción por frentes: montaje a 1,5 camiones/día en terminal con Célula Beta, plataforma elástica en AKS/IoT Hub para 450 viajes/día, y turnos 4x4 en Torre 24/7 con peak de 37 FTE nivelado en T-15; enlazado desde Formulario T-6).*}{S1 §1.4 (01-empresa/contenido.tex) y Formulario T-6 (Informe-3/d1/). Estado: Incorporada en D1}
\observacion[sub=E]{EMP-02}{EMP-02. El gerente general es también jefe de proyecto y representante legal con dedicación de 100 \% entre M1 y M20. Falta explicar la conducción de la empresa en ese período. (rev., pp. 5)}{Documentar la cobertura de funciones corporativas y las delegaciones durante la dedicación contractual. Conciliar cargos y disponibilidad con la nivelación de recursos. *(Implementado en D1: incorporada en S1 §1.3 la subsección «Régimen de delegación corporativa y continuidad institucional» con Tabla 1-delegacion-funciones asignando subrogancias y poderes durante M1–M20 mientras Alejandro Hermosilla mantiene dedicación contractual exclusiva).*}{S1 §1.3 (01-empresa/contenido.tex), Tabla 1-delegacion-funciones. Estado: Incorporada en D1}
\observacion[sub=E]{EMP-03}{EMP-03. Se vinculan entidades certificadoras y auditoras reales con expedientes ficticios. (rev., pp. 5, 47)}{Usar entidades inequívocamente ficticias para los antecedentes simulados y armonizar sus referencias, identificadores y respaldos en todo el expediente. *(Implementado en D1: reemplazadas todas las entidades reales por ficticias únicas: Bureau Veritas \$\rightarrow\$ Andes Certificaciones Técnicas SpA, exp. ACT-EXP-2026-CL-8921; Nexia \$\rightarrow\$ Austral Auditores Consultores SpA en S1 §§1.3, 1.4, 1.5 y Anexo 1.A).*}{S1 §§1.3–1.5 y Anexo 1.A (01-empresa/anexo-1a.tex). Estado: Incorporada en D1}
\observacion[sub=E]{EMP-04}{EMP-04. Las alianzas exclusivas se remiten a un capítulo y un formulario ausentes del Informe 2. (rev., pp. 5)}{Incorporar las alianzas del contrato en el capítulo de equipo y subcontrataciones y adjuntar T-8 para la oferta final. Conciliar alcance y remisiones con las alianzas presentadas en S1. *(Implementado en D1: reformulado S1 §1.6 delimitando que las alianzas corporativas [Azure CSP Tier 1, Quectel, Technoton, talleres Ruta 5] respaldan solvencia técnica, mientras que las cartas de patrocinio y subcontratación contractual se radican en Sobre N.º 1 y se resumen en Formulario T-8, eliminando dependencias circulares).*}{S1 §1.6 (01-empresa/contenido.tex). Estado: Incorporada en D1}
\observacion[sub=E]{NEC-01}{NEC-01. La estimación de 338,15 viajes diarios de temporada alta y la cota de 450 viajes diarios dependen de un supuesto de 210 viajes diarios en valle. (rev., pp. 7–8)}{Calibrar con registros del sistema de gestión de 2013 y conciliar la distribución con los 96.000 viajes anuales. Si los registros siguen indisponibles, documentar procedencia, sensibilidad y margen del escenario, y recalcular la carga adoptada donde se utilice. *(Implementado en D1: formalizado el canon de 96.000 viajes/año en S2 §2.1 y la ecuación de balance en §2.2.3; caracterizada la cota de 450 viajes/día estrictamente como cota de estrés de diseño ingenieril [+33,1 \% sobre media alta] para absorber ráfagas sin pérdida en Event Hubs, programando la calibración empírica con TMS 2013 en Hito H1 [Actividad A04] y diferenciando viajes/día de las 380/570 sesiones de usuario en Torre de Control evaluadas en CP-PERF-01).*}{S2 §§2.1 y 2.2.3 (02-problema/contenido.tex). Estado: Incorporada en D1}
\observacion[sub=E]{NEC-02}{NEC-02. El diagnóstico presenta los umbrales de continuidad como compromisos propios y usa la expresión «análisis pericial». (rev., pp. 7–8)}{S2 atribuye a las Bases la disponibilidad, el RTO, el RPO y la autonomía de la solución. Remite al diseño, la operación y las pruebas correspondientes. Se retiró la expresión «análisis pericial».}{S2 §§2.2.2 y 2.5.1, ff. 13 y 26–27. Estado: Incorporada}
\observacion[sub=E]{NEC-03}{NEC-03. La figura de cadena de valor queda aislada por espacios en blanco. (rev., pp. 7)}{La Figura 2.1 se redibujó en formato vectorial y se integró entre el párrafo que la anuncia y la lectura por fases. Se comprobó el texto a 9,5 pt en el PDF.}{S2 §2.2, Figura 2.1, f. 10. Estado: Incorporada}
\grupoObservaciones{Alcance y cumplimiento técnico}
\observacion[sub=E]{ALC-01}{ALC-01. SA-01 fija febrero de 2027 como inicio y el calendario depende de esa fecha para la ventana de M16. (rev., pp. 11–12, 46–47)}{Separar los meses contractuales del calendario civil y evaluar inicios alternativos contra las ventanas de montaje y puesta en producción. Si la combinación resulta incompatible, elevar la consulta formal sin cambiar unilateralmente los hitos ni supeditar la oferta a SA-01.}{S3 supuestos y calendario, S7 §§7.2–7.3, T-14/T-15/T-18 y E-25. Estado: Sin cierre acreditado}
\observacion[sub=E]{ALC-02}{ALC-02. La reversión se resume en una afirmación de que no se pierde información. (rev., pp. 11)}{Identificar disparadores, puntos de restauración, tratamiento de transacciones y conciliación de evidencia. Remitir desde S3 al procedimiento operativo y a su prueba. *(Implementado en D1: enriquecido el caso de prueba CP-SYS-20 en T-17 articulando el procedimiento de reversión de 6 pasos de Formulario T-18: detección por mesa de ayuda en 15 min, decisión con contraparte en 30 min, desvío de tráfico en API Gateway a versión N-1 en 15 min, replay de buffer de 72 h a 20 min/unidad y conciliación transaccional con ERP en 24 h sin pérdida de datos).*}{Formulario T-17 (09-calidad/formularios/catalogo-T-17.tex), caso CP-SYS-20. Estado: Incorporada en D1}
\observacion[sub=E]{ALC-03}{ALC-03. El T-12 contiene 42 requerimientos propios y usa los RT como referencias, sin declarar su cumplimiento individual. (rev., pp. 11–12, 33, 46–47)}{Inventariar cada requisito del caso y cada RT del capítulo 3 transversal. Declarar cumplimiento individual, componente, sección y folio de evidencia. Conciliar el inventario con T-17 y justificar expresamente cualquier no aplicabilidad.}{T-12, S3 trazabilidad y T-17. Estado: Sin cierre acreditado}
\observacion[sub=E]{ALC-04}{ALC-04. El T-12 declara «Sí» uniforme, incluso en RF-003, RNF-005 y RF-014, que dependen de autorizaciones o pruebas futuras. (rev., pp. 11–12, 46–47)}{Distinguir compromiso, evidencia disponible y verificación futura. Declarar el cumplimiento parcial cuando corresponda, explicitar la dependencia y definir cómo se resuelve, revisando su compatibilidad con el Art. 53 para los requisitos obligatorios.}{T-12 fichas RF-003, RNF-005 y RF-014, S3 supuestos, T-16 y T-17. Estado: Sin cierre acreditado}
\observacion[sub=E]{ALC-05}{ALC-05. La numeración de las tablas 3.1–3.7 del T-12 coincide con la del S3. (rev., pp. 11–12, 47)}{Usar un prefijo propio de T-12 y actualizar cada remisión en ambos documentos.}{T-12 y S3. Estado: Sin cierre acreditado}
\observacion[sub=E]{ALC-06}{ALC-06. El cuerpo del S3 se refiere a versiones previas del informe. (rev., pp. 11–12)}{Redactar el cuerpo como propuesta autónoma y concentrar la historia de las correcciones en este anexo.}{S3 alcance, implementación e implantación, y anexo del Art. 46. Estado: Sin cierre acreditado}
\grupoObservaciones{Arquitectura y datos}
\observacion[sub=E]{ARQ-01}{ARQ-01. Flutter, React y Next.js no tienen versiones y soporte definidos. La versión Go declarada carece de respaldo verificable para la fecha de oferta. (rev., pp. 13–14, 43)}{Fijar versiones verificadas para la fecha de oferta. Incorporar fuente de soporte, fecha de fin o política fechada y plan de actualización durante el contrato. Alinear el inventario de componentes con el T-11.}{S4 §4.1, Tabla de tecnologías, y T-11. Estado: Sin cierre acreditado}
\observacion[sub=E]{ARQ-02}{ARQ-02. Los ocho ADR llevan una misma fecha y carecen de estado por decisión. (rev., pp. 14–15)}{Registrar por ADR su fecha efectiva, estado, alternativa seleccionada, descartes y criterio, con evidencia de la decisión y su vigencia.}{S4 §4.1, registro ADR y Tabla 4.2. Estado: Sin cierre acreditado}
\observacion[sub=E]{ARQ-03}{ARQ-03. La mesa de ayuda y la dotación operativa se remiten al S10 sin dimensionamiento. (rev., pp. 15–16, 18, 44)}{Derivar contactos por franja, cargas y turnos de atención, dotación, suplencias y escalamiento. Conciliar capacidad y cobertura 24×7 con el plan de operación y sus SLA.}{S4 §4.2, S10–S11 y T-15. Estado: Sin cierre acreditado}
\observacion[sub=E]{ARQ-04}{ARQ-04. El alcance de ExpressRoute hasta Chile Central queda sujeto a confirmación del proveedor. (rev., pp. 17–18, 29, 31, 46–47)}{Obtener confirmación documentada del circuito y probar su conmutación. Si la ruta ofertada no resulta viable, acreditar una conectividad conforme a las Bases y registrar la exposición y su respuesta en T-16.}{S4 §4.2, conectividad, T-11, S8/T-16 y T-17. Estado: Sin cierre acreditado}
\observacion[sub=E]{ARQ-05}{ARQ-05. El DR de datos personales en Brazil South carece de vía habilitante documentada. R-24 depende de una actuación del mandante antes de H2 y afecta H7. (rev., pp. 17–20, 30–31, 45–47)}{Documentar la vía aplicable y su respaldo. Definir un sitio o mecanismo admisible con evidencia de RTO y RPO si la habilitación no se obtiene. Vincular la decisión de H2, la contingencia y el efecto sobre H7, sin asumir autorización del mandante.}{S4 §4.2 DR, S5 §5.1.6 soberanía y tratamiento, S8/T-16 R-24, S7/T-18 y T-17. Estado: Sin cierre acreditado}
\observacion[sub=E]{ARQ-06}{ARQ-06. La tabla del T-11 se identifica con el mismo número que la tabla de tecnologías de S4 y remite las referencias al capítulo. (rev., pp. 17–18, 47)}{Asignar numeración propia al formulario y localizar sus fuentes de especificación y soporte, manteniendo remisiones explícitas a S4 cuando procedan.}{T-11 y S4, referencias y tabla de tecnologías. Estado: Sin cierre acreditado}
\observacion[sub=E]{ARQ-07}{ARQ-07. Persisten brechas de cobertura real, parque de terceros, conductor sin rastro, consultas sin respuesta y autonomía de la solución satelital referida en la revisión. (rev., pp. 18)}{Identificar para cada brecha el dato que falta, la prueba o consulta, el plazo de resolución y el diseño de contingencia. Separar la autonomía de la solución ofertada de antecedentes de versiones anteriores.}{S4 §§4.1–4.2, S3 supuestos y decisiones, T-16 y T-17. Estado: Sin cierre acreditado}
\observacion[sub=E]{DAT-01}{DAT-01. El diccionario de datos se incorporó dentro del archivo del S5. (rev., pp. 20, 47)}{Publicar el diccionario como anexo separado y conservar en el cuerpo los modelos, el análisis y la remisión al anexo.}{S5 y AUDIT-Subdocumento5-Anexos. Estado: Sin cierre acreditado}
\observacion[sub=E]{DAT-02}{DAT-02. Las tablas 5.2, 5.3, 5.5 y 5.6 contienen prosa explicativa en las celdas. (rev., pp. 20, 47)}{Trasladar las reglas y sus excepciones al texto o al diccionario. Reservar las celdas para atributos comparables y códigos.}{S5, Tablas 5.2, 5.3, 5.5 y 5.6. Estado: Sin cierre acreditado}
\observacion[sub=E]{DAT-03}{DAT-03. La declaración de S5 deja una página de continuación con encabezado y sin filas. (rev., pp. 20, 47)}{Eliminar el salto vacío y revisar el cierre, los encabezados y los folios en el PDF recompilado.}{S5, Declaración de uso de IA y folio final. Estado: Sin cierre acreditado}
\observacion[sub=E]{DAT-04}{DAT-04. El 99,9 \% mensual de servicios críticos en S5 debe cotejarse con el SLA de S4. (rev., pp. 20)}{Conciliar meta, servicios cubiertos, ventana y fórmula de cálculo entre datos, arquitectura, operación y sus pruebas.}{S5 Tabla 5.4, S4 SLA, S10 y T-17. Estado: Sin cierre acreditado}
\grupoObservaciones{Metodologías}
\observacion[sub=E]{MET-01}{MET-01. El plan de comunicaciones cubre cinco grupos y no articula los trece actores de S2 ni define destinatario, frecuencia o canal. (rev., pp. 21, 23)}{Desarrollar una matriz por interesado con información, responsable, destinatario, frecuencia, canal y escalamiento. Conciliarla con los trece actores del diagnóstico.}{S6 §6.1, T-9 y S2 análisis de interesados. Estado: Sin cierre acreditado}
\observacion[sub=E]{MET-02}{MET-02. Falta identificar quién prioriza y acepta incrementos por Curimón y quién lo reemplaza. (rev., pp. 21, 23)}{Definir con la contraparte el rol titular y suplente, atribuciones y tiempos de decisión. Distinguir aceptación de incremento de aceptación contractual y vincular las revisiones con H4 y H9.}{S6 §§6.1–6.2, T-9/T-10 y S7 plan de entregas. Estado: Sin cierre acreditado}
\observacion[sub=E]{MET-03}{MET-03. T-9 y T-10 se limitan a remisiones, con poca información utilizable por sí mismos. (rev., pp. 21–22)}{Adjuntar a T-9 la información de gestión de S6 §6.1 y a T-10 la de desarrollo de §6.2, conforme a sus instrucciones. Identificar contenido y evidencias sin depender de remisiones genéricas.}{T-9 y T-10, con S6 §§6.1–6.2. Estado: Sin cierre acreditado}
\observacion[sub=E]{MET-04}{MET-04. Los ocho sprints de A09 y cinco de A20 no tienen un plan de entregas por contexto o requerimiento. Sólo se detalla G13 en el sprint 7. (rev., pp. 22–23, 26–27, 47)}{Asignar a cada sprint requisitos, incremento, dependencias, capacidad neta, prueba, usuario que revisa y criterio de salida. Conciliar G13 y los cierres de H4/H9 con el T-15 y T-17.}{S6 §6.2, S7 §§7.1–7.3, T-10, T-12, T-14/T-15 y T-17. Estado: Sin cierre acreditado}
\observacion[sub=E]{MET-05}{MET-05. La deuda técnica se descuenta en sprints, pero carece de paquete y horas durante la operación Kanban. (rev., pp. 22–23)}{Incluir trabajo de deuda técnica durante M21–M56 en el plan de operación. Dimensionar esfuerzo y capacidad sin confundir la regla de capacidad Scrum con la asignación Kanban.}{S6 §6.2, S7 §7.2, T-14/T-15 y S11. Estado: Sin cierre acreditado}
\observacion[sub=E]{MET-06}{MET-06. La viñeta de ramas temáticas de GitFlow cita normas ISO que no sustentan esa práctica. (rev., pp. 23)}{Retirar las citas improcedentes y respaldar el flujo de ramas con la fuente técnica pertinente. Citar ISO/IEC 25012 donde se aplique a calidad de datos.}{S6 §6.2, control de versiones y referencias. Estado: Sin cierre acreditado}
\observacion[sub=E]{MET-07}{MET-07. Las tablas 6.1, 6.2, 6.3 y 6.6 contienen prosa. La Tabla 6.2 utiliza una descripción arquitectónica que no coincide con S4. (rev., pp. 23)}{Trasladar al texto las reglas extensas y armonizar la descripción de contenedores y contextos con los ADR del S4.}{S6, Tablas 6.1, 6.2, 6.3 y 6.6, y S4 §4.1. Estado: Sin cierre acreditado}
\grupoObservaciones{EDT, cronograma e implantación}
\observacion[sub=E]{PLA-01}{PLA-01. El primer nivel de la EDT se justifica genéricamente, divide el software por etapa y carece de una agrupación intermedia. (rev., pp. 24, 27, 47)}{Justificar cada rama por entregable y agrupar el software por capacidad o contexto, con subgrupos y paquetes verificables. Representar la etapa como atributo de planificación. La vista revisada aún conserva ramas por etapa.}{S7 §7.1 y T-14. Evidencia parcial en S7 revisado, ff. 7 y 23. Estado: Parcial}
\observacion[sub=E]{PLA-02}{PLA-02. Las láminas EDT repetían el diccionario en vez de mostrar la jerarquía. (rev., pp. 24, 27)}{El S7 revisado muestra el árbol general y cinco vistas de descomposición por rama. El diccionario conserva los criterios y responsables en T-14.}{S7 revisado §7.1, f. 7, Figuras 7.3–7.8, ff. 23–28, y T-14. Estado: Incorporada}
\observacion[sub=E]{PLA-03}{PLA-03. Los cinco paquetes de innovación carecían de meses concretos en el T-14 y de actividad identificable en la red. (rev., pp. 24, 27, 37)}{El T-14 revisado fecha los cinco paquetes y la Gantt de S7 muestra sus ventanas. Integrar en T-15 el trabajo y las dependencias que afectan otros paquetes, y cotejar las ventanas con la versión final de T-19.}{T-14 revisado f. 7 y S7 revisado f. 40. Destino adicional T-15 y T-19. Estado: Parcial}
\observacion[sub=E]{PLA-04}{PLA-04. El esfuerzo se calcula desde dotaciones mensuales y no por tarea y paquete. (rev., pp. 25, 27, 47)}{Estimar horas por entregable, tarea, rol y cantidad de trabajo. Conciliar el total con la curva por etapa y la curva separada de operación, verificando picos y solapes.}{S7 §7.2 y T-15, detalle de esfuerzo y curvas. Estado: Sin cierre acreditado}
\observacion[sub=E]{PLA-05}{PLA-05. A09 y A20 agrupan toda la construcción de cada etapa sin secuenciar los contextos dependientes. (rev., pp. 25, 27, 47)}{Descomponer ambas actividades y enlazar vigencias y jornada con asignación, documento y costo, según sus dependencias reales. Recalcular holguras, ruta crítica y efecto en hitos. La red revisada aún conserva las actividades agregadas.}{S7 §§7.2–7.3 y T-15 revisado, Tabla 7.7, ff. 8–10. Destino adicional T-14 y T-16. Estado: Sin cierre acreditado}
\observacion[sub=E]{PLA-06}{PLA-06. La migración histórica y el estudio de costo por ruta no tienen precedencias ni holguras en la red. (rev., pp. 25)}{Incorporar 7.2 y 2.5 a la red con predecesoras, sucesoras, responsables y criterio de término. Recalcular margen y fecha de entrega a las áreas usuarias.}{S7 §§7.2–7.3, T-14 paquetes 7.2 y 2.5 y T-15. Estado: Sin cierre acreditado}
\observacion[sub=E]{PLA-07}{PLA-07. La Gantt por paquete no muestra dependencias ni la ruta crítica de manera conciliable con T-15. (rev., pp. 25, 27, 47)}{La vista CPM revisada incorpora ruta crítica y enlaces principales para las actividades agregadas. Vincular los paquetes con la red en una Gantt que muestre también sus dependencias y su contribución a la ruta crítica.}{S7 revisado §7.3.2 f. 19, vista CPM f. 33 y vistas de paquetes ff. 34–43. Destino adicional T-15. Estado: Parcial}
\observacion[sub=E]{PLA-08}{PLA-08. La ruta crítica encadena levantamiento, arquitectura, infraestructura, construcción y certificación con escaso paralelismo. (rev., pp. 25)}{Evaluar inicios parciales por capacidad y por ambiente. Mantener como predecesoras las dependencias técnicas demostrables y justificar las actividades que deban seguir en serie.}{S7 §§7.2–7.3 y T-15, red CPM y supuestos. Estado: Sin cierre acreditado}
\observacion[sub=E]{PLA-09}{PLA-09. Las láminas Gantt de pocas filas y las 42 páginas gráficas fragmentaban la lectura. (rev., pp. 25, 27)}{El S7 revisado concentra la Gantt en un panorama, una vista CPM y diez vistas temáticas. La EDT se presenta aparte en el mismo bloque de anexos gráficos.}{S7 revisado §7.3.2 y Figuras 7.12–7.23, ff. 32–43. Estado: Incorporada}
\observacion[sub=E]{PLA-10}{PLA-10. La primera marcha blanca coincide con la temporada de fruta sin declarar la disponibilidad de usuarios de torre y terminales como dependencia. (rev., pp. 26)}{Acordar ventanas, suplencias y dedicación con el mandante y registrar la contingencia si falta participación. Conservar los hitos y las condiciones de salida del Art. 17.3.}{S7 §7.3, T-18, S8/T-16 y S9 coordinación de UAT. Estado: Sin cierre acreditado}
\observacion[sub=E]{PLA-11}{PLA-11. A18 combina integración telemática en M13–M15 con montaje en M16–M18, pese a la restricción estacional. (rev., pp. 26–27)}{El S7 revisado distingue las ventanas en el texto. Separar ambas actividades también en la red, con recursos, precedencias y capacidad de montaje por terminal.}{S7 revisado §7.3.2 f. 19 y T-15 revisado, Tabla 7.7, f. 9. Destino adicional T-14/T-18. Estado: Parcial}
\observacion[sub=E]{PLA-12}{PLA-12. El piloto A11 coincide con meses de cierres de Los Libertadores. La holgura se justifica con un cierre individual de hasta 12 días. (rev., pp. 26)}{Analizar cierres sucesivos, su coincidencia con el piloto y su consumo de holgura. Definir disparador y alternativa de ruta o reprogramación que preserve la aceptación del piloto.}{S7 §7.3, T-15 actividades A11 y A18, S8/T-16 y T-18. Estado: Sin cierre acreditado}
\observacion[sub=E]{PLA-13}{PLA-13. El plan de montaje no incorpora las 210 balizas de semirremolque, las 21 de reposición ni los lectores de portería declarados para la innovación 3. (rev., pp. 27, 47)}{Asignar paquete, sitios, actividad, personal y rendimiento a balizas y lectores. Distinguir unidades de montaje de stock de reposición y conciliar cantidades y ventanas con T-11 y T-19.}{S7 §§7.1–7.3, T-14 paquete 4.2, T-15, T-18, T-11 y T-19 innovación 3. Estado: Sin cierre acreditado}
\grupoObservaciones{Riesgos}
\observacion[sub=E]{RIE-01}{RIE-01. La mitad del registro del cliente agrupa riesgos internos de audIT. R-06 trata como incertidumbre un retiro anunciado y ya resuelto por diseño. (rev., pp. 28, 31, 47)}{Clasificar por origen, exposición residual y actuación requerida del mandante. Mantener controles propios en arquitectura, planificación y registro interno. Reclasificar R-06 como antecedente de la decisión de caché y conservar visibles los riesgos contractuales que sigan abiertos.}{S8 §§8.1–8.3, T-16, S4 ADR de caché y registro interno de audIT. Estado: Sin cierre acreditado}
\observacion[sub=E]{RIE-02}{RIE-02. La convivencia y el retiro del sistema de 2013 no aparecen en el registro de riesgos. (rev., pp. 29, 31, 47)}{Registrar la exposición por función en M16, M21 y M24, con disparadores, conciliación, reversión, responsables y efecto sobre los hitos.}{S8/T-16, S3 estrategia de sustitución, S7 y T-18. Estado: Sin cierre acreditado}
\observacion[sub=E]{RIE-03}{RIE-03. Faltan en el registro los riesgos de adopción de las innovaciones 1 y 2 y su riesgo residual. (rev., pp. 29, 31)}{Evaluar utilidad percibida, errores de estimación de llegada e incompatibilidad del kit. Vincular cada evento, disparador y respuesta con la ficha T-19 y recalcular el residual.}{S8/T-16 y S13/T-19, innovaciones 1 y 2. Estado: Sin cierre acreditado}
\observacion[sub=E]{RIE-04}{RIE-04. El árbol de fallas de asignación a conductor no apto carece de cuantificación. (rev., pp. 29, 31)}{Definir antes del paralelo de M13–M14 cómo se cuentan los eventos, su exposición y las dependencias. Cuantificar con los datos obtenidos y declarar incertidumbre y límites de la estimación.}{S8 §8.2, Figura 8.3, T-16 y plan de prueba en paralelo T-18/T-17. Estado: Sin cierre acreditado}
\observacion[sub=E]{RIE-05}{RIE-05. El análisis costo-beneficio de respuestas no presenta una razón de retorno. (rev., pp. 30)}{Calcular costo incremental, pérdida esperada evitada y sensibilidad de cada respuesta material. Mantener los montos y la razón monetaria en el análisis económico y dejar el criterio técnico en S8.}{S8 §8.3, T-16 y análisis económico del Sobre 3. Estado: Sin cierre acreditado}
\observacion[sub=E]{RIE-06}{RIE-06. R-17 se acepta con seguimiento, aunque la simulación reporta 21,6 \% de atraso de H1 y 16,1 \% de H2. (rev., pp. 30–31, 47)}{Definir umbrales de alerta, escalamiento, recuperación y contingencia para ambos hitos. Recalcular el residual con la red actualizada sin cambiar la aceptación contractual de los hitos.}{S8 §8.3 y T-16 R-17, S7 y T-15. Estado: Sin cierre acreditado}
\observacion[sub=E]{RIE-07}{RIE-07. Las tablas 8.1–8.3 se repiten entre S8 y T-16 y las remisiones son ambiguas. (rev., pp. 31)}{Usar numeración propia del T-16 y reconciliar las referencias a contingencias, disparadores y respuestas.}{S8 §8.3 y T-16, tablas de registro y respuesta. Estado: Sin cierre acreditado}
\grupoObservaciones{Calidad y aceptación}
\observacion[sub=E]{CAL-01}{CAL-01. El calendario de ensayos y las puertas de calidad se muestran sólo en tablas de detalle. (rev., pp. 32)}{S9 incorpora una vista temporal de campañas y puertas alineada con S7/T-15. La figura y su explicación distinguen controles recurrentes, ambientes, entradas, evidencias de salida e hitos de aceptación.}{S9 §9.3.2, Figura 9.2 y explicación, ff. 13–14. Estado: Incorporada}
\observacion[sub=E]{CAL-02}{CAL-02. CP-HW-08 alude a polvo minero y CP-HW-06 duplica el ensayo de alarma SOS de CP-PERF-06. (rev., pp. 32, 34)}{CP-HW-08 identifica el ensayo IP67 sobre el conjunto instalado en transporte terrestre. CP-HW-06 verifica entrega SOS por evento con cobertura real. CP-PERF-06 mide el P99 bajo degradación reproducible. Se conservaron los 120 identificadores del catálogo.}{T-17, CP-PERF-06 f. 76, CP-HW-06 f. 94 y CP-HW-08 ff. 95–96. Estado: Incorporada}
\observacion[sub=E]{CAL-03}{CAL-03. Los 29 criterios de aceptación del capítulo 18 del caso no se vinculan a casos de prueba ni a actas. (rev., pp. 33–34, 47)}{Trazar cada criterio a requisito, prueba, umbral, evidencia, responsable e hito. Distinguir la cobertura de criterios de la cobertura de los 42 requisitos propios. *(Implementado y corregido en D1: creada y depurada la matriz de trazabilidad completa de los 29 criterios en T-17 [Tabla 9-rtm-29]; corregido el Criterio 28 desacoplándolo del ensayo SOS CP-HW-06 y reasignándolo a CP-UNIT-18, CP-SYS-07 y CP-UAT-07; e individualizados los responsables de verificación de audIT y Curimón en las 29 filas).*}{Formulario T-17 (09-calidad/formularios/trazabilidad-T-17.tex), Tabla 9-rtm-29. Estado: Incorporada en D1}
\observacion[sub=E]{CAL-04}{CAL-04. La UAT de Etapa 1 coincide con el comienzo de la temporada de fruta, sin disponibilidad acordada de usuarios de terminal y despacho. (rev., pp. 34)}{Programar ventanas UAT, cupos por rol, suplencias y escalamiento. Registrar la consulta y el plan de contingencia si la disponibilidad de usuarios no se confirma.}{S9 §§9.2–9.3, T-13/T-17, S7/T-18 y T-16. Estado: Sin cierre acreditado}
\observacion[sub=E]{CAL-05}{CAL-05. El S9 ubica A12 en M10–M12 y el T-13 la expresa como M12/H5. (rev., pp. 34)}{S9 y T-13 distinguen la ejecución de A12 en M10–M12 de la certificación H5 en M12. Se cotejó la ventana con la actividad A12 de T-15.}{S9 Tabla 9.5, f. 14. T-13, apartado 5, f. 8. Estado: Incorporada}
\observacion[sub=E]{CAL-06}{CAL-06. El T-17 incluye un protocolo general y carece de protocolos por cada hito y del producto final. (rev., pp. 34, 47)}{Preparar una ficha por H1–H12 y por producto final con entregable, criterio objetivo, casos, evidencia, responsable, plazos, observaciones y acta de conformidad. *(Implementado y corregido en D1: formalizadas las 13 fichas completas de protocolo de aceptación H1 a H12 y Producto Final en protocolos-hitos-T-17.tex; corregidos los porcentajes y la base de cobro calculada directamente sobre el valor total de la fase de implementación [60\% Etapa 1, 40\% Etapa 2] conforme a E-25, separando los 36 cobros mensuales de operación M21–M56).*}{Formulario T-17 (09-calidad/formularios/protocolos-hitos-T-17.tex, T-17.tex) y S9 §9.3.3. Estado: Incorporada en D1}
\grupoObservaciones{Innovaciones y valorización}
\observacion[sub=E]{INN-01}{INN-01. La programación por oportunidad de patio tiene TRL 2 y metas relativas a una línea base sin medir. (rev., pp. 35–37, 45)}{Definir línea base y método de medición en piloto, umbrales observables y fecha de decisión. Si la solución no demuestra beneficio o madurez suficiente, definir y validar un sustituto que cumpla el tipo exigido.}{S13 y T-19 innovación 2, T-14/T-15 y análisis económico. Estado: Sin cierre acreditado}
\observacion[sub=E]{INN-02}{INN-02. La innovación 3 presenta como novedad el kilometraje real de semirremolques, que apoya el criterio de aceptación 26. (rev., pp. 36–37)}{Clasificar el kilometraje como cumplimiento del criterio 26 y delimitar la innovación a verificación de acople y registro térmico adicional. Separar sus pruebas, indicadores y beneficios.}{S13/T-19 innovación 3, T-12 y T-17 criterio 26. Estado: Sin cierre acreditado}
\observacion[sub=E]{INN-03}{INN-03. La innovación 5 presenta el catálogo de paradores del criterio 28 como parte del mérito innovador. (rev., pp. 35–37, 45)}{Ubicar el catálogo en el alcance obligatorio. Delimitar la innovación a anticipación circadiana y al modelo no punitivo, con indicadores y pruebas separados. *(Implementado en D1: actualizada la Ficha 5 en Formulario T-19 delimitando que el catálogo base de paraderos seguros pertenece al alcance obligatorio del Criterio 28 [Art. 25 bis], restringiendo el mérito innovador estrictamente a la modelación circadiana adaptativa y protocolo de notificación preventiva no punitiva).*}{Formulario T-19 (13-innovaciones/formularios/T-19.tex), Ficha 5. Estado: Incorporada en D1}
\observacion[sub=E]{INN-04}{INN-04. Las cinco fichas omiten el riesgo residual. La percepción de vigilancia en la innovación 5 se describe con probabilidad media a alta e impacto alto. (rev., pp. 36–37)}{Indicar riesgo inicial, tratamiento, verificación y residual en cada ficha. Para innovación 5 definir una medida adicional con criterio de activación y verificar el efecto de la notificación acústica y el tratamiento no punitivo. *(Implementado en D1: incorporada en Ficha 5 de T-19 la matriz de riesgos formal con severidad inherente 3x4 [Alta], gobernanza de retroalimentación no punitiva, auditoría de privacidad y riesgo residual mitigado 1x3 [Baja]).*}{Formulario T-19 (13-innovaciones/formularios/T-19.tex), Ficha 5. Estado: Incorporada en D1}
\observacion[sub=E]{INN-05}{INN-05. Las innovaciones 2, 3 y 5 carecen de base y magnitud para sus beneficios, aunque el informe remite la valorización al Sobre 3. (rev., pp. 36)}{La programación de patio se evaluará mediante visitas y minutos de preparación evitados. En la innovación 3 se medirán por separado la verificación de acople y el registro térmico adicional. En la innovación 5 se medirán los desvíos evitados atribuibles a la recomendación y el consumo asociado, y se contrastará la siniestralidad con registros históricos. La cuantificación de los beneficios queda pendiente. La inversión y la operación se registran por separado del beneficio de Curimón.}{S13 y T-19, fichas 2, 3 y 5. Documento de costos frente a venta, sección de valorización de las innovaciones, y planilla económica. Estado: Parcial}
\end{tablaObservaciones}
